\SetPageTheme{story}
\PageHeader{Calculatorul nu stie sa citeasca}{Povestea cu care incepem}{Deschidere}

\begin{missionbox}[Inceput narativ]
  Cand un om citeste un mesaj, vede imediat cuvinte, sens si intentie. Calculatorul nu vede nimic din toate acestea. Pentru el, un mesaj este doar o succesiune ordonata de simboluri care trebuie analizate exact.

  Daca acel mesaj contine o parola, un cod de verificare sau o informatie care trebuie protejata, calculatorul are nevoie de o metoda foarte stricta de lucru. El trebuie sa stie unde incepe informatia, unde se termina, cum o imparte in bucati si ce compara cu ce.

  Aici se afla una dintre ideile fundamentale ale confidentialitatii digitale: protejarea informatiei presupune mai intai control foarte bun asupra formei ei.
\end{missionbox}

\begin{conceptbox}[Ce invatam de fapt]
  Problema palindromului este doar un pretext bun. In spatele ei invatam:
  \begin{itemize}
    \item cum descompune calculatorul informatia;
    \item cum urmareste ordinea elementelor;
    \item cum compara doua structuri;
    \item cum construieste o verificare pas cu pas.
  \end{itemize}
\end{conceptbox}

\begin{guidebox}[Activare]
  Scrie in doua-trei propozitii de ce crezi ca un calculator are nevoie de reguli stricte pentru a prelucra mesaje, parole sau coduri.
  \AnswerSpace{3}
\end{guidebox}

\newpage
\SetPageTheme{theory}
\PageHeader{Recapitulare Minima}{Ce luam cu noi in lectii}{Fundatie}

\begin{conceptbox}[Instrumente minime]
  In aceasta editie folosim foarte putine idei de baza:
  \begin{itemize}
    \item variabile pentru a retine starea;
    \item operatori pentru a extrage informatie;
    \item \texttt{daca} pentru decizii simple;
    \item \texttt{cat timp} pentru repetitie;
    \item tabele de trace pentru intelegere.
  \end{itemize}
\end{conceptbox}

\begin{examplebox}[Variabilele cele mai importante]
\begin{lstlisting}[style=codequest]
int n, copie, uc, ogl, cnt, suma;
\end{lstlisting}
\end{examplebox}

\begin{guidebox}[Roluri]
  Completeaza:
  \begin{itemize}
    \item \texttt{n} pastreaza \hrulefill
    \item \texttt{copie} pastreaza \hrulefill
    \item \texttt{uc} inseamna \hrulefill
    \item \texttt{ogl} va deveni \hrulefill
  \end{itemize}
\end{guidebox}
